\section*{Changelog}
\addcontentsline{toc}{section}{Changelog}

\subsubsection*{2026-11-11}
\begin{itemize}
  \item Added TWIC Mode~2 Active Card Authentication as a PIV Auto profile, with a capability flag,
    certificate-derived FASC-N, and a proposed 127-byte challenge and RSA/ECC prehash construction.
    Added NexGen V10 AID recognition and clarified unknown minor-version handling. Updated the
    Appendix~C ACA example and documented Mode~4 composition using commanded fingerprint
    verification.
  \item Specified maximum card response-length requests as \code{0xFF} (short) and \code{0xFFFF}
    (extended), with shorter nonzero lengths permitted. Documented TSA's interpretation of ISO
    response-length encoding, its conflict with TWIC Appendix~E, and potential compatibility issues
    with other federally issued credentials.
  \item Renamed the OSDP channel requirement to Secure Channel~3.0. Required protected TPK loading,
    biometric-object reads, biometric matching, and match replies with deployment-specific keys,
    per-fragment checks, and cleanup on channel loss or replacement to protect the ACU--PD link
    against line tapping.
  \item Required biometric signature validation and identity binding for PIV and TWIC. Allowed
    maintained enrollment-time validation, with the reply hash identifying the verified reference
    record.
  \item Made PIVMODE flag \code{0x02} require SM for PUT DATA and kept SELECT ordinary. Added
    protected and ordinary provisioning examples, with completion tied to the card's result.
  \item Enabled contact and contactless NEXGEN TWIC free read/write E-Stickers in the selected TWIC
    application. Defined native TLV encoding, accepted \code{0x6282} for complete reads, and added
    write/read examples and a captured write/read/erase sequence.
  \item Distinguished TWIC fingerprint verification, Reference Picture Verification, and PACS-held
    alternate-biometric matching, with FASC-N and transaction binding under 33~CFR~101.105
    and~101.520(c).
  \item Specified TWIC AES-128-ECB decryption, padding and reference checks. Defined Result
    \code{0x07} for TPK/decryption failure, \code{0x0A} for an absent reference, and \code{0x0B} for
    an unusable reference.
  \item Required multiple-card rejection, PIVERROR \code{0x102A}, and a one-second RF removal
    interval between TWIC presentations.
  \item Required CS2 and CS7 for credential SM preparation and defined biometric working-data and
    TPK zeroization. Specified the SHA-256-derived card identifier and coordinate-only KDF and
    key-confirmation inputs, with the uncompressed-point prefix retained in the card command.
  \item Defined Function Code~15 capability bits for legacy SPE, explicit PIV SPE, and ACU-armed
    Auto Execute. Simplified LOADTPK to a 16-byte key for one biometric operation in the current
    presentation, with unconditional key cleanup.
  \item Aligned PIVMODE's Bluetooth interface bit with Function Code~13 at \code{0x10}, and reserved
    \code{0x08} in both global and per-entry masks.
  \item Updated TWIC Part~4 to October 2026, corrected registration and operational citations, and
    added the 2026 EC384 technical advisories. Referenced 33~CFR~101.520 for authentication and
    validity checks. Corrected digest byte-size notation in the autonomous-PIV table.
\end{itemize}

\subsubsection*{2026-10-14}
\begin{itemize}
  \item Added per-profile key references and certificate-based algorithm discovery so the same PIV
    Auto configuration supports RSA and ECC cards.
  \item Added a resume option after pairing or PIN entry, defined counter handling and the challenge
    format, and updated the appendix with PIV and Legacy/NextGen TWIC examples.
\end{itemize}

\subsubsection*{2026-10-07}
\begin{itemize}
  \item Added the Invalid TPK result for TWIC Privacy Key failures and set the TPK field length to
    16 bytes.
  \item Clarified PIV/TWIC application selection and added a Data Model byte to PIVSELECT.
  \item Updated the Legacy/NextGen fingerprint and facial workflows and the BIOMATCH hash
    description.
  \item Clarified CRAUTH cleanup so each new request starts a fresh transaction while the previous
    reply remains available for retransmission.
\end{itemize}

\subsubsection*{2026-09-16}
\begin{itemize}
  \item Renamed \texttt{osdp\_PIVDATA}/\texttt{osdp\_PIVDATAR} to
    \texttt{osdp\_PIVGETDATA}/\texttt{osdp\_PIVGETDATAR}, retaining command code \code{0xA3}. Added
    identifier and tag lengths, two-byte offsets and requested lengths, and whole-object or unique
    direct-child selection with the original tag and length bytes preserved. Defined OSDP 3.0
    multi-part delivery and an all-zero first reply for empty success. Command data length
    distinguishes the new request from the fixed-length OSDP 2.2 form.
  \item Added Format Marker \code{0xFF} to the reassembled \texttt{osdp\_CRAUTH} request, retaining
    command code \code{0xA5}. Multi-part lengths and offsets cover the marker, algorithm, key
    reference, and challenge. OSDP 3.0 ACUs use the new CRAUTH and PIVGETDATA forms with supporting
    PDs; PDs may additionally support the OSDP 2.2 forms, with format-specific response handling.
  \item Renamed \texttt{osdp\_PDERROR} to \texttt{osdp\_PIVERROR} and added GET DATA selection and
    transfer errors. The new CRAUTH and PIVGETDATA forms return one terminal PIVERROR on failure,
    replacing CRAUTH's empty failure reply and optional follow-up error. Renumbered tentative
    PIV-specific errors consecutively from \code{0x1020} through \code{0x1029}; credential removal
    is now \code{0x1026}.
  \item Separated credential SM setup, VCI pairing, and PIN verification. Clarified automatic
    preparation over contact and contactless, PIV Auto failure reporting, and ACU-directed recovery.
    Authenticated PIN and pairing rejections preserve SM; SM processing and integrity failures end
    the session.
  \item Added \texttt{osdp\_CARDSTATUS} Type \code{0x02} diagnostics containing the SM outcome and
    credential CVC hash. Every CARDSTATUSR now includes Status Type after reply code \code{0x85}:
    \code{0x01} for standard status and pairing completion, \code{0x02} before the multi-part fields
    in each diagnostic fragment. Extended \texttt{osdp\_PIVSTATUS} with SW1/SW2 when a verified
    credential failure status is available; the status word follows the signed response in the
    reassembled payload.
  \item Defined cached-PIN reuse without extending the original deadline, silent idle expiration,
    and completion of requests affected by expiration. Bounded PD timeout defaults to 1--255
    seconds. Reassigned SPE flags \code{0x01} for retention across removal and \code{0x02} for Auto
    Execute; added result \code{0x0B} for credential communication failure or timeout, requiring
    preparation to restart after an unconfirmed PIN VERIFY. Applied the same restart rule to pairing
    VERIFY timeouts, reported through existing PIVERROR \code{0x0002} (smart-card timeout).
  \item Added workflow diagrams and worked packet examples for selection, SM/VCI recovery, PIN-cache
    timing, and multi-part replies. Standardized table sizes and protocol values on hexadecimal
    notation.
\end{itemize}

\subsubsection*{2026-09-02}
\begin{itemize}
  \item Made proposal builds self-contained with the SIL-licensed Carlito typeface in regular, bold,
    italic, and bold-italic faces. Documented the XeLaTeX and \texttt{latexmk} requirements.
  \item Completed the Enhanced PIV message-name cleanup by applying the compact OSDP naming
    convention consistently to \code{osdp_PIVLOADTPK}, \code{osdp_PIVXMITPAIRING},
    \code{osdp_PIVVCILOADTA}, \code{osdp_PIVVCITASTATUS}, and \code{osdp_PIVVCITASTATUSR}; command
    and reply codes, payloads, and behavior are unchanged.
  \item Standardized multi-part messages on the generic OSDP \emph{MpSizeTotal}, \emph{MpOffset},
    and \emph{MpFragmentSize} field names and moved the multi-part fields ahead of the reassembled
    \code{osdp_PIVMODE} and \code{osdp_PIVVCILOADTA} data.
  \item The current \code{osdp_PIVVCILOADTA} definition supersedes the anchor-only multi-part sizing
    described in the 2026-05-30 entry: its total now includes the 12-byte metadata block.
  \item Updated the PIV Auto simulator, generated fixtures, proposal PDF, and redline PDF to
    exercise and publish the same framing and naming.
\end{itemize}

\subsubsection*{2026-08-12}
\begin{itemize}
  \item Proposed clear, mode-specific SPE behavior for status, PIN acquisition and verification,
    cached-PIN verification, cache/abort, and attempts-remaining queries. The proposal also defines
    finite default timeouts, cache and PIN-verified lifecycles, and ACK/poll completion behavior.
  \item Proposed a one-byte \texttt{osdp\_PIVSPER} reply whose value is interpreted by the requested
    mode. The OSDP~2.2 \texttt{osdp\_PR01SPER} reply remains unchanged.
  \item Standardized Enhanced PIV names, including renaming the autonomous PIV Auto reply to
    \texttt{osdp\_PIVSTATUS}. Corrected tentative code conflicts by assigning
    \texttt{osdp\_PIVSELECT} to \code{0xB0} and \texttt{osdp\_PDERROR} to \code{0x8A}.
\end{itemize}

\subsubsection*{2026-06-18}
\begin{itemize}
  \item Added \texttt{osdp\_PIVSELECT} as an OSDP-native direct AID selection command using
    tentative command code \code{0xA7}; clarified its ISO/IEC~7816 \texttt{SELECT by DF name}
    mapping, selected-AID status update, and application-state clearing rules.
  \item Removed stale \texttt{osdp\_CRAUTHR} \texttt{osdp\_PDERROR} wording so codes
    \code{0x0006}/\code{0x0007} retain their common meanings for secure-channel and multi-part
    errors.
  \item Clarified that VCI pairing codes are transmitted from the ACU to the PD only over OSDP
    Secure Channel~3.0 and are used by the PD to establish VCI secure messaging with the credential.
  \item Tightened and clarified verbiage by replacing imprecise ``legacy'' terminology throughout
    normative text with neutral wording such as unchanged, OSDP~2.2-format, and poll-based; retained
    ``legacy'' only for withdrawn algorithms such as RSA-1024 test material.
\end{itemize}

\subsubsection*{2026-05-30}
\begin{itemize}
  \item Added NIST-derived VCI and PIV APDU fixture coverage, GSA-backed demonstration keys, and
    reproducible worked examples for PIV Auto and VCI validation. RSA-1024 material is retained only
    as informative legacy/non-current test-vector material.
  \item Clarified PIV Auto cryptographic inputs and response handling, including KMAC-derived
    challenge construction, RSA and ECC \texttt{GENERAL AUTHENTICATE} templates, nonzero UUID
    requirements, and complete returned \code{0x7C} Dynamic Authentication Templates in
    \texttt{osdp\_PIVSTATUS}.
  \item Aligned related command semantics for capability advertising, \texttt{osdp\_CARDSTATUS}
    errors, selective \texttt{osdp\_PIVDATA} replies, \texttt{osdp\_PIVSPE} status/abort modes, and
    VCI trust-anchor inventory.
  \item Reworked Appendix C and the references section for standards review: examples are explicitly
    informative, vector provenance is documented, source fixtures are framed as public test
    material, and normative references now include the needed APDU, KMAC, X.509, and
    requirement-language specifications.
  \item Added supporting scripts and tests for deterministic simulation, fixture replay, VCI
    import/validation, Appendix C generation, YubiKey material loading, live APDU capture, and
    reproducible PDF/table builds.
\end{itemize}

\subsubsection*{2026-05-27}
\begin{itemize}
  \item Version renamed from OSDP~2.3 to OSDP~3.0 throughout.
  \item Removed \texttt{osdp\_PIVSPE} flag \code{0x01} (retry continuation); each
    \texttt{osdp\_PIVSPE} command now initiates exactly one PIN-entry session and
    \texttt{osdp\_PIVSPER} terminates the session unconditionally, including on PIN mismatch. ACUs
    wishing to allow additional attempts \textbf{MAY} re-issue \texttt{osdp\_PIVSPE}, using the
    \emph{Attempts Remaining} count from the preceding \texttt{osdp\_PIVSPER} to guide retry policy.
    Bit \code{0x01} is now reserved.
  \item Aligned \texttt{osdp\_PIVSPER} and \texttt{osdp\_PIVBIOMATCHR} result codes for credential
    removal, changed \texttt{osdp\_PIVBIOMATCHR} \code{0x03} to PIN verification required, reserved
    \texttt{osdp\_PIVBIOMATCHR} \code{0x07}, clarified biometric assurance limits, and specified
    that only off-card biometric match success returns the 32-byte SHA-256 hash of tag \code{0xBC}.
  \item Clarified the Function Code~14 biometric-support reference for \texttt{osdp\_PIVLOADTPK} and
    added CAC as PIV mode data model \code{0x03}.
  \item Added PIV Auto KMAC challenge derivation, required series and sequence counters, fixed-width
    supplemental entropy, and the \texttt{osdp\_PIVSTATUS} poll reply for reporting credential
    identity and signed challenge responses. PIV Auto is configured through \texttt{osdp\_PIVMODE}
    and is restricted to Secure Channel~3.0 sessions that are not using factory keys.
  \item Replaced the SP~800-73 Intermediate CVC anchor model with an OSDP-defined VCI trust anchor
    record (tag \code{0x7F50}) carrying a Profile Identifier (\code{0x80} = \code{0x01}), an
    eight-byte Issuer Identifier Number (\code{0x42}), and a DER-encoded SubjectPublicKeyInfo
    (\code{0x7F49}). Maximum conforming record size is 576~bytes (\code{0x0240}), derived from the
    RSA-4096 key case.
  \item Added normative IIN derivation rule: the IIN in \texttt{osdp\_PIVVCILOADTA} and in the trust
    anchor record \textbf{SHALL} be the leftmost eight bytes of the subjectKeyIdentifier from the
    content-signing X.509 certificate.
  \item Tightened \texttt{osdp\_PIVVCILOADTA} multi-part encoding: Anchor Length and MpSizeTotal
    ranges corrected to \code{0x0000}-\code{0x0240}; added normative constraint that
    \texttt{MpOffset}~+~\texttt{MpFragmentSize} \textbf{SHALL NOT} exceed \texttt{Anchor Length}; PD
    \textbf{SHALL} verify the reassembled byte string is exactly \emph{Anchor Length} bytes and a
    single \code{0x7F50} record before committing.
  \item Documented variable-length PD storage accounting: PDs \textbf{SHALL} account for anchors by
    their actual \emph{Anchor Length}, not a fixed maximum; replacement is atomic (prior anchor
    remains if replacement fails); delete releases actual stored bytes.
  \item Specified the two-tier PD verification path during VCI establishment: direct (trust anchor
    IIN matches secure messaging CVC IIN, anchor SPKI verifies CVC signature) and intermediate (no
    direct match triggers retrieval of PIV data object \code{0x1017}/\texttt{5FC122}; Intermediate
    CVC Subject Identifier must equal CVC IIN, Role \code{0x5F4C} must be \code{0x12}, anchor IIN
    must match Intermediate CVC IIN, anchor SPKI verifies Intermediate CVC, Intermediate CVC
    verifies secure messaging CVC). PDs \textbf{SHALL NOT} perform X.509 path validation or
    revocation for VCI trust anchors.
  \item Added normative ACU post-VCI validation requirements: after \texttt{osdp\_CARDSTATUSR}
    reports VCI established, the ACU \textbf{SHALL} obtain the content-signing certificate (from
    cache or via \texttt{osdp\_PIVDATA} for \code{0x1017}/inner tag \code{0x70}), verify SKI matches
    \emph{VCI Anchor IIN}, verify SPKI matches the stored anchor record, verify PIV or PIV-I
    content-signing EKU OID, verify validity period (expired certificate \textbf{SHALL} cause
    session rejection and anchor unload), and perform X.509 path validation and revocation per
    deployment policy.
  \item Restored \emph{VCI Trust Anchor ID} to \texttt{osdp\_CARDSTATUSR} with normative language;
    added \emph{VCI CVC IIN} and \emph{VCI Anchor IIN} fields (each 8~bytes) so the ACU receives the
    IIN pair needed to distinguish direct and intermediate verification paths.
  \item Updated \texttt{osdp\_PIVVCITASTATUSR}: clarified that \emph{Available Bytes} reflects
    variable-length accounting and that the per-record IIN is the anchor selector from the stored
    trust anchor record, not necessarily the secure messaging CVC IIN reported in
    \texttt{osdp\_CARDSTATUSR}.
  \item Added VCI security preconditions: all VCI-related commands \textbf{SHALL} be rejected unless
    OSDP Secure Channel~3.0 is active and not using factory/default keys; failures reported via
    \texttt{osdp\_PDERROR} \code{0x0006}.
\end{itemize}

\subsubsection*{2026-01-21}
\begin{itemize}
  \item Error-handling cleanup and code assignments: \texttt{osdp\_PIVPUTDATA} now \textbf{SHALL}
    report failures via \texttt{osdp\_PDERROR} (status word optional) so writes can’t silently
    succeed.
  \item \texttt{osdp\_PIVSPE} adds result code \code{0x09} in \texttt{osdp\_PIVSPER} for “VCI not
    established.”
  \item PDERROR is credential/global only; PD capability gaps stay \texttt{osdp\_NAK}.
  \item OSDP~2.2-format replies (e.g., \texttt{osdp\_CRAUTHR}, \texttt{osdp\_PIVDATAR}) remain
    unchanged, with optional follow-up \texttt{osdp\_PDERROR} detail if the PD has it.
  \item Documented tentative command/reply codes for the new PIV set: commands \code{0xA6} (VCI
    trust-anchor status) through \code{0xAF} (VCI trust-anchor load), and replies
    \code{0x85}-\code{0x89} for the new responses.
\end{itemize}

\subsubsection*{2026-01-07}
\begin{itemize}
  \item Added a PIV Mode interface-mask table aligned to Function Code~13 bits.
  \item Clarified CRAUTHR's full GENERAL AUTHENTICATE response payload, with status words carried
    through PDERROR.
  \item Clarified CARDSTATUS and CARDSTATUSR behavior: solicited status after PD preparation, with
    an error on credential removal.
  \item Remapped PIV error codes to the \code{0x10xx} block and added PIV APDU references.
  \item Aligned PIVDATA selection of Discovery (\code{0x7E}) and BIT (\code{0x7F61}) subtags with
    SP~800-73-4.
  \item Clarified PDERROR's single-error reply and PIVXMITPAIRING submission before VCI
    establishment when the credential requires pairing.
\end{itemize}

\subsubsection*{2025-12-10}
\begin{itemize}
  \item Consolidated application-level errors on \texttt{osdp\_PDERROR} (two-byte codes, optional
    SW1\textbar SW2 for APDU failures, multi-error support)
  \item moved the PIV error set to \code{0x0100}-\code{0x01FF} and updated PIV commands to use
    \texttt{osdp\_PDERROR} (transport faults remain \texttt{osdp\_NAK})
  \item made the Global Interface Mask in \texttt{osdp\_PIVMODE} a required byte (no TLV tag) with
    an updated field layout
  \item clarified that \texttt{osdp\_CARDSTATUSR} is solicited-only after PD preparation with
    poll-based \texttt{osdp\_RAW} as the aggregated card-present indication
  \item and noted the deprecation of \texttt{osdp\_GENAUTH} with a recommendation to migrate to
    \texttt{osdp\_CRAUTH}.
\end{itemize}

\subsubsection*{2025-12-04}
\begin{itemize}
  \item Clarified that \texttt{osdp\_CRAUTHR} returns the full \texttt{GENERAL AUTHENTICATE} data
    including the Dynamic Authentication Template (excluding only the ISO/IEC~7816 status word).
\end{itemize}

\subsubsection*{2025-11-12}
\begin{itemize}
  \item Added PIV Biometric Match request/reply and tables (one-bit type values only)
  \item moved Enhanced PIV NAKs to a 0x20+ base (added 0x27 Credential Not Present/Removed, 0x28
    Biometric Type Not Supported; clarified 0x23 Credential Security Status Not Satisfied)
  \item aligned CRAUTHR error handling to reference the enhanced NAK set
  \item relocated Enhanced PIV (0x04) and PIV Auto (0x08) indicators to Function Code 12 and cleaned
    Function Code 13 to interface-only bits (removed non-interface flags; removed PIV Auto from
    Function Code 14)
  \item updated CARDSTATUS gating to Function Code 12 bit 0x04
  \item added a PIV Mode Flags table and clarified PIV Mode AID encoding
  \item enabled PDF TOC hyperlinks and improved table/caption pagination.
\end{itemize}

\subsubsection*{2025-10-30}
\begin{itemize}
  \item Clarified CRAUTH command/reply behaviour (including RSA, AES, and ECC support), overhauled
    PIV Mode configuration with TLV-based application and interface controls (including PIV Auto),
    added card status request/reply messages, expanded PIV data retrieval semantics, introduced PIV
    Put Data, PIV Secure PIN Entry, TWIC privacy-key load, VCI pairing transmit, and VCI
    trust-anchor load/status commands, and documented the enhanced error-code set (including VCI/PIV
    Auto conditions).
\end{itemize}

\subsubsection*{2025-09-17}
\begin{itemize}
  \item Major PIV Mode and SPE revisions.
\end{itemize}

\subsubsection*{2025-06-25}
\begin{itemize}
  \item Added multi-AID usage and VCI support notes.
\end{itemize}

\subsubsection*{2025-05-27}
\begin{itemize}
  \item New function code definitions for 3.0 biometric handling.
\end{itemize}
